\section{A linear bound on distinct normal component vectors}
\label{sec:component-profile-bound}

Weighted orbit extraction produces a compressed inventory even when the
number of connected components is enormous. For an actual embedded normal
surface, its fully coalesced inventory satisfies a stronger, linear bound.
We prove it from normal-coordinate uniqueness and elementary normal
neighborhoods. This is a finiteness refinement for the present analysis;
we make no claim of literature-wide priority. Classical normal-surface
finiteness is discussed by King~\cite[Theorem 2]{king2003}, who attributes
the underlying result to Haken and Kneser and notes that stronger bounds
were already known.

Let $\mathcal T$ be a finite triangulation of a compact $3$-manifold $M$,
possibly with boundary. Tetrahedral faces may be identified within or
between tetrahedra, but the quotient must be a manifold: every vertex
link is a connected sphere or disc. Ideal vertices, singular vertex links,
and immersed surfaces are outside this statement. Write $t$ for the
number of tetrahedra and $v$ for the number of quotient vertices.
Orientability of $M$ is not required.

For a properly embedded normal surface $F$, its standard vector
$\mathbf x(F)\in\mathbb Z_{\ge0}^{7t}$ records four triangle counts and
three quadrilateral counts in each tetrahedron. These counts satisfy the
standard matching equations and the quadrilateral constraints. Standard
vectors determine normal surfaces up to ambient normal isotopy, including
their connected component multisets~\cite[Lemma 2.5 and Theorem 2.10]{burton2009}.
Let $Q(F)$ be the set of nonzero quadrilateral coordinates and put
$q=|Q(F)|\le t$. Count a component vector only once, regardless of how
many components have that vector, and denote the resulting number by $H$.

\begin{proposition}[Component-type bound]
\label{prop:component-type-bound}
Under these hypotheses,
\[
                     H\le v+2q.
\]
If every component of $F$ is two-sided, then $H\le v+q$.
Neither conclusion assumes incompressibility, least weight, or
fundamentality of $F$.
\end{proposition}

\begin{lemma}[Independence of disjoint two-sided types]
\label{lem:disjoint-normal-independence}
Mutually disjoint, connected, two-sided normal surfaces with pairwise
distinct standard vectors have linearly independent standard vectors.
\end{lemma}

\begin{proof}
Write the vectors as $\mathbf y_1,\ldots,\mathbf y_h$. Since their entries
are integers, a nontrivial real linear dependence implies a nontrivial
rational dependence: row reduction of the corresponding integer matrix
computes its kernel over $\mathbb Q$. Clearing denominators gives an
integer dependence. The vectors are nonzero and coordinatewise
nonnegative, so the nonzero coefficients cannot all have the same sign.
Consequently a dependence would give
\[
 \sum_{i\in P}a_i\mathbf y_i
       =\sum_{j\in R}b_j\mathbf y_j,
 \qquad a_i,b_j\in\mathbb Z_{>0},\quad P\cap R=\varnothing,
\]
with both index sets nonempty.

Choose mutually disjoint, sufficiently small normal product neighborhoods
of the surfaces. Two-sidedness permits any positive number of parallel
normal copies inside each such neighborhood, each with the same standard
vector as the original. The two sums therefore represent actual disjoint
unions having precisely the displayed component multisets. Equality of
their full vectors gives a normal isotopy between the unions. Such an
isotopy sends each connected component to a connected component and
preserves its normal-disc counts. Thus some $\mathbf y_i$ on the left
equals a $\mathbf y_j$ on the right, contradicting distinctness.
\end{proof}

\begin{proof}[Proof of Proposition~\ref{prop:component-type-bound}]
Let $V_Q\subseteq\mathbb R^{7t}$ be the real solution space of the
standard matching equations with all quadrilateral coordinates outside
$Q(F)$ set to zero. Project $V_Q$ to its $q$ permitted quadrilateral
coordinates. Its kernel has dimension exactly $v$.

Indeed, when every quadrilateral coordinate vanishes, face matching
identifies the triangle coefficients at the corners joined through that
face. The corner triangles around each quotient vertex have connected
adjacency graph, because its link is a connected surface. Their
coefficients must therefore agree. Values at different quotient vertices
are independent. The resulting basis consists of the vertex-link vectors
$\boldsymbol\ell_1,\ldots,\boldsymbol\ell_v$, each with one triangle at
every corner of its vertex and no other discs. Rank--nullity gives
\begin{equation}\label{eq:normal-support-dimension}
                  \dim V_Q\le v+q.
\end{equation}
In particular, every connected triangle-only normal surface is one
vertex link.

Choose one representative of each component type of $F$ having a nonzero
quadrilateral coordinate. Retain a two-sided representative unchanged.
Replace a one-sided representative $S$ by the horizontal frontier of a
small normal interval neighborhood. The frontier is the connected double
cover associated to the nontrivial normal line bundle of $S$; it is
two-sided and has standard vector $2\mathbf x(S)$. For a surface with
boundary, the vertical boundary of its interval neighborhood lies in
$\partial M$, and the horizontal frontier is properly embedded.

All these neighborhoods can be chosen mutually disjoint and disjoint
from retained representatives. Coalesce equal resulting vectors, and let
$r$ be their number. Adjoin all $v$ vertex links, taking them sufficiently
close to the triangulation vertices to avoid every chosen surface. The
resulting $r+v$ surfaces are mutually disjoint, connected and two-sided.
Their vectors are distinct: the first $r$ have nonzero quadrilateral
part, whereas vertex links do not. They lie in $V_Q$ and are independent
by Lemma~\ref{lem:disjoint-normal-independence}. Thus
\[
                r+v\le\dim V_Q\le v+q,
                \qquad r\le q.
\]

Each resulting nonvertex vector $\mathbf z$ has at most two original
preimages: a two-sided type with vector $\mathbf z$, and a one-sided
type with vector $\mathbf z/2$. Two different one-sided vectors cannot
have the same double. Hence there are at most $2r\le2q$ original
nonvertex types, and at most $v$ triangle-only types. This proves the
first bound. If all components are two-sided, the replacement map is
injective and the same argument gives $H\le v+q$.
\end{proof}

\subsection{A sharp one-tetrahedron example}

The factor two accounts for a real geometric possibility: a one-sided
component and its connected two-sided double can coexist. The fixture
\texttt{layered\_torus(1)} makes this explicit. Label the tetrahedron
vertices $0,1,2,3$, with face $i$ opposite vertex $i$, and identify face
$0$ with face $1$ by the permutation
$(0\,1\,2\,3)$; faces $2$ and $3$ remain boundary faces. This is a
one-vertex solid torus. Use coordinate order
\[
 (T_0,T_1,T_2,T_3;Q_{01|23},Q_{02|13},Q_{03|12}).
\]
Its vertex-link disc and a normal M\"obius band have vectors
\[
 \boldsymbol\ell=(1,1,1,1;0,0,0),\qquad
 \mathbf m=(0,0,0,0;0,1,0).
\]
For the single quadrilateral of $\mathbf m$, the corners on edges
$03,23,12$ become one vertex, while its corner on $01$ is another.
There are three edge cells and one face cell, so its Euler characteristic
is zero. The two unpaired sides form one boundary circle. The quotient
is therefore a M\"obius band. Its normal frontier is an annulus with
vector $2\mathbf m$.

Choose this band, its small frontier, and the vertex link disjointly.
Their union has vector
\[
 \boldsymbol\ell+3\mathbf m=(1,1,1,1;0,3,0)
\]
and exactly the three distinct component vectors
$\boldsymbol\ell,\mathbf m,2\mathbf m$. Thus $H=3=v+2q$ with $v=q=1$.
The union with vector $\boldsymbol\ell+2\mathbf m$ also attains the
two-sided bound $H=2=v+q$. The supplied native inventory and independent
Regina decomposition agree on these explicit components; the full fixture
and audit record are retained in
\texttt{results/normal\_inventory\_audit.json}.

\subsection{Consequences for encoded output}

Suppose the supplied coordinates have at most $B\ge1$ bits. Every
component coordinate is bounded by its source coordinate. A component
multiplicity is at most the total number of normal discs, hence needs
$O(B+\log(t+1))$ bits. Since $v\le4t$ and $q\le t$, the coalesced
normal inventory has
\[
             O\bigl(t^2(B+\log(t+1))\bigr)
\]
bits. The number of components itself remains unbounded for fixed
triangulation; it is their number of distinct full vectors that is linear.

AHT already establish polynomial-time recovery of component data from
binary normal coordinates~\cite[Corollary 17]{aht2006}. The present bound
sharpens the final output accounting and permits any polynomial-time
predicate on one connected vector to be evaluated on every type.
It does not enumerate candidate source surfaces. Nor does equality of
the weighted sum of proposed component vectors certify decomposition:
normal sums can reconnect surfaces, so actual component membership still
comes from the checked orbit calculation.
